\section{Trivializing maps}

Somewhat surprisingly, trivializing maps can, to some extent,
be constructed expli\-cit\-ly in the pure gauge theory.
The construction is explained in this section, assuming
that the gauge action $S(U)$ is
a sum of Wilson loops (plaquettes, rectangles, etc.).


\subsection{Trivializing flows}

If the generator $Z_t(U)$ of the flow (3.2) is such that
\begin{equation}
  \int_0^t\rmd s\,\sum_{x,\mu}
  \bigl\{\du^a_{x,\mu}[Z_s(U)]^a(x,\mu)\bigr\}_{U=U_s}
  =tS(U_t)+C_t,
\end{equation}
where $C_t$ may depend on $t$ but not on the fields,
the associated integrated transformations satisfy
\begin{equation}
  S(\transt(V))-\ln\det\transts(V)=(1-t)S(\transt(V))-C_t.
\end{equation}
In particular, the transformation at $t=1$
is then a trivializing map.

Equation (4.1) is a rather implicit condition on the generator of
the flow. However, when differentiated with respect to
$t$, it assumes a more tractable form,
\begin{equation}
  \sum_{x,\mu}\bigl\{
  \du^a_{x,\mu}[Z_t(U)]^a(x,\mu)
  -t\du^a_{x,\mu}S(U)[Z_t(U)]^a(x,\mu)\bigr\}=
  S(U)+\dot{C}_t,
\end{equation}
which involves the generator at time $t$ only.
Note that the differential condition (4.3)
and the flow equation (3.2) imply eq.~(4.1),
i.e.\ it suffices to find a generator $Z_t(U)$ that satisfies
eq.~(4.3).


\subsection{Existence of trivializing flows}

Equation (4.3) is an inhomogeneous linear partial differential
equation for the generator $Z_t(U)$.
Since it is a scalar equation,
one expects that there are many solutions.
In the following, the solution will be obtained in the form
\begin{equation}
  [Z_t(U)]^a(x,\mu)=-\du^a_{x,\mu}\Sflow_t(U),
\end{equation}
where the action $\Sflow_t(U)$ is to be determined.

When inserted in eq.~(4.3), the ansatz (4.4) leads to the
Laplace equation
\begin{align}
  &\Lop\Sflow_t=S+\dot{C}_t,
  \\[1.0ex]
  &\Lop=
  \sum_{x,\mu}\bigl\{
  -\du^a_{x,\mu}\du^a_{x,\mu}
  +t\left(\du^a_{x,\mu}S\right)\du^a_{x,\mu}\bigr\}
\end{align}
(for simplicity, the argument $U$ is now often omitted).
The operator $\Lop$ is elliptic and symmetric with respect to
the scalar product
\begin{equation}
   (\phi,\psi)=\int\rmD[U]\,\rme^{-tS(U)}\phi(U)^*\psi(U).
\end{equation}
$\Lop$ has therefore a complete set of differentiable
eigenfunctions and a purely discrete spectrum
with no accumulation points (see ref.~\cite{Gilkey}, sect.~1.6,
for example).
Moreover, since
\begin{equation}
   (\phi,\Lop\phi)=\sum_{x,\mu}\bigl(\du^a_{x,\mu}\phi,\du^a_{x,\mu}\phi
   \bigr)\geq0,
\end{equation}
the function $\phi(U)=1$ is the only zero mode of $\Lop$
and all other eigenfunctions have eigenvalues separated
from the origin by a strictly positive spectral gap.

Now if one chooses $C_t$ to be such that
\begin{equation}
   \dot{C}_t=-(1,S)/(1,1),
\end{equation}
the zero-mode component is removed from
the right-hand side of eq.~(4.5) and
\begin{equation}
   \Sflow_t=\Lop^{-1}(S+\dot{C}_t)
\end{equation}
is then a well-defined expression that solves the equation.
The differentiability of $\Sflow_t$ with respect to $t$ and $U$
essentially follows from the ellipticity of $\Lop$
(appendix~E).
A constructive proof of the existence of
trivializing flows has thus been given.


\subsection{Expansion in powers of $t$}

The solution (4.10) is well defined but still quite
implicit since it involves the inverse of an operator
acting on functions of the gauge field. In this and
the next subsection,
the solution is worked out analytically
in powers of $t$.

When the series
\begin{equation}
  \Sflow_t=\sum_{k=0}^{\infty}t^k\Sflow^{(k)}
\end{equation}
is inserted in eq.~(4.5),
the matching of the terms of a given order in $t$ leads to the
recursion
\begin{align}
  \Lap\Sflow^{(0)}&=S+\dot{C}^{(0)},
  \\[1.0ex]
  \Lap\Sflow^{(k)}&
  =-\sum_{x,\mu}\du^a_{x,\mu}S\,\du^a_{x,\mu}\Sflow^{(k-1)}
  +\dot{C}^{(k)},\qquad
  k=1,2,\ldots,
\end{align}
for the actions $\Sflow^{(k)}$.
The Laplacian $\Lap$ coincides with the colour-electric
part of the Hamilton operator in lattice gauge theory
in $4+1$ dimensions.
In particular, its eigenfunctions are products
of $\Group$ representation
functions of the link variables.
Sums of
Wilson loops and products of Wilson loops, for example,
are eigenfunctions of $\Lap$
or can easily (algebraically) be
decomposed into eigenfunctions.

The solution of the recursion,
\begin{align}
  \Sflow^{(0)}&=\Lap^{-1}S,
  \\[1.0ex]
  \Sflow^{(k)}&=-\Lap^{-1}
  \sum_{x,\mu}\du^a_{x,\mu}S\,\du^a_{x,\mu}
  \Sflow^{(k-1)},\qquad k=1,2,\ldots,
\end{align}
is thus obtained in the form of sums of
Wilson loops and products of Wilson loops.
Note that,
as already mentioned in subsect.~4.2,
the constant function is the only zero mode of $\Lap$.
Since the smallest non-zero eigenvalue
of $\Lap$ is $\frac{4}{3}$,
the right-hand sides of
eqs.~(4.14),(4.15) are therefore unambiguously determined
up to an irrelevant additive constant.


\subsection{Calculation of\/ $\Sflow^{(0)}$ and\/ $\Sflow^{(1)}$
                in the Wilson theory}

For illustration, the first two terms of the
series (4.11)
are now worked out explicitly for the plaquette action~\cite{Wilson}
\begin{equation}
  \Sw(U)=-\frac{1}{6}\beta\,\Wl_0(U)
\end{equation}
(where $\beta=6/g_0^2$ denotes the inverse gauge coupling).
A short calculation,
using the completeness relation (A.5),
shows that the leading term is
\begin{equation}
  \Sflow^{(0)}=-\frac{1}{32}\beta\Wl_0=\frac{3}{16}\Sw.
\end{equation}
To this order and up to
a rescaling of the time parameter $t$, the trivializing
flow in the Wilson theory thus coincides with the Wilson flow.

\begin{figure}[htbp]
  \centering
  \includegraphics[width=8.5cm]{images/loops.pdf}
  \caption{%
Classes of loops and pairs of loops contributing to $\Sflow^{(1)}$
in the Wilson theory.
The loops $5-7$
reside on a single plaquette of the lattice.
All other loops occupy two
plaquettes which can lie in a plane or be at right angles
in three dimensions.
}
  \label{fig:loops}
\end{figure}

The expression to be worked out at the next order is
\begin{equation}
  \Sflow^{(1)}=
  -\frac{1}{192}\beta^2\Lap^{-1}\sum_{x,\mu}\du^a_{x,\mu}\Wl_0\,
  \du^a_{x,\mu}\Wl_0.
\end{equation}
The Wilson loops and products of Wilson loops
that can occur at this point derive from the
contractions of two plaquette loops with a common link.
Altogether there are then seven classes $\mathcal{C}_i$, $i=1,\ldots,7$, of
loops and pairs of loops to consider (see fig.~2).
By summing the traces of the associated Wilson loops,
each class $\mathcal{C}_i$ defines an action
\begin{align}
  \Wl_i&=\sum_{C\in\mathcal{C}_i}\tr\bigl\{U(C)\bigr\}
  \quad\text{if}\quad i=1,2,5,
  \\[1.0ex]
  \Wl_i&=\sum_{\{C,C'\}\in\mathcal{C}_i}\tr\bigl\{U(C)\bigr\}\tr\bigl\{U(C')\bigr\}
  \quad\text{if}\quad i=3,4,6,7,
\end{align}
where $U(C)$ denotes the ordered product of the link variables
along the loop $C$. The sums in these equations extend over all possible
positions of the loops on the lattice. Loops with opposite
orientation are considered to be different and are both included
in the sums.

Using the identity (A.5) again, some algebra now yields
\begin{multline}
  \sum_{x,\mu}\du^a_{x,\mu}\Wl_0\,\du^a_{x,\mu}\Wl_0
  =\Wl_1-\Wl_2-\frac{1}{3}\Wl_3+\frac{1}{3}\Wl_4\\
  -2\Wl_5+\frac{2}{3}\Wl_6-\frac{4}{3}\Wl_7
\end{multline}
up to an additive constant. Furthermore,
\begin{align}
  \Lap\Wl_1&=8\Wl_1,
  \\[0.6ex]
  \Lap\Wl_2&=\frac{31}{3}\Wl_2+\Wl_4,
  \\[0.6ex]
  \Lap\Wl_3&=11\Wl_3-\Wl_1,
  \\[0.6ex]
  \Lap\Wl_4&=\frac{31}{3}\Wl_4+\Wl_2,
  \\[0.6ex]
  \Lap\Wl_5&=\frac{28}{3}\Wl_5+4\Wl_6,
  \\[0.6ex]
  \Lap\Wl_6&=\frac{28}{3}\Wl_6+4\Wl_5,
  \\[0.6ex]
  \Lap\Wl_7&=12\Wl_7+\text{constant}.
\end{align}
In the subspace of these functions,
the operator $\Lap$ can be easily inverted
and the result
\begin{multline}
  \Sflow^{(1)}=\frac{1}{192}\beta^2\Bigl\{
  -\frac{4}{33}\Wl_1
  +\frac{12}{119}\Wl_2
  +\frac{1}{33}\Wl_3
  -\frac{5}{119}\Wl_4
  +\frac{3}{10}\Wl_5\\
  -\frac{1}{5}\Wl_6
  +\frac{1}{9}\Wl_7\Bigr\}
\end{multline}
is thus obtained.
Note that the smallness of the numerical coefficients in this formula
is balanced, to some extent,
by the number of loops per lattice point
in the classes $\mathcal{C}_i$ (which are equal to
$120$, $12$ and $6$ for $i=1,\ldots,4$, $i=5,6$ and $i=7$ respectively).


\subsection{Miscellaneous remarks}

\noindent
(a)~{\itshape Higher orders}.
The actions $\Sflow^{(k)}$, $k\geq2$, can
be computed algebraically
following the steps taken in the previous subsection.
Since all loops generated by contracting
a plaquette loop with the loops at order $k-1$
must be considered, the work required for
the calculation
is however rapidly growing with $k$.

\smallskip
\noindent
(b)~{\itshape Locality and convergence}. The series (4.11) is an
expansion in local terms whose footprint on the lattice
increases proportionally to the order $k$.
At the values of $t$, where the expansion converges,
the action $\Sflow_t$ is then guaranteed to be local as well.

The norm estimates in appendix E imply a lower bound on
the convergence radius of the series, but this bound
is rather poor and vanishes in the infinite-volume limit.
It seems nevertheless plausible that
the series has a non-zero convergence radius
in this limit,
because the inverse of the operator $\Lop$
in eq.~(4.10) is likely to remain bounded
in a complex neighbourhood of $t=0$.
An analysis that
takes the locality properties of $\Lop$
into account will however be required to show this.

\smallskip
\noindent
(c)~{\itshape Truncation of the expansion}.
If all terms in the series (4.11) of order $k\geq n$ are dropped,
one obtains an approximately trivializing flow that
satisfies eq.~(4.2) up to an additive correction
of order $t^{n+1}$.
An at least partial cancellation of the action is
achieved in this case as long as $t$ is
sufficiently small for the correction
to be strongly suppressed.

\smallskip
\noindent
(d)~{\itshape Smoothing property}.
The Wilson flow satisfies
\begin{equation}
  {\rmd\over\rmd t}\Sw(U_t)=-\frac{3}{16}\sum_{x,\mu}
  \bigl\{\du^a_{x,\mu}\Sw(U)
  \du^a_{x,\mu}\Sw(U)\bigr\}_{U=U_t}\leq 0
\end{equation}
and therefore lowers the Wilson action as $t$ increases.
To leading order in $t$, the trivializing flow
constructed in this section for the Wilson theory thus
has a smoothing effect on the gauge field.
On the other hand, if the flow is followed in the reverse
direction, the gauge field tends to become rougher.

\smallskip
\noindent
(e)~{\itshape Topological charge sectors}.
In lattice QCD, the topological (instanton)
sectors are not a property of the field manifold alone, but
are expected to emerge dynamically when the continuum limit is approached.
The fact that trivializing maps completely ``straighten out''
the sectors is therefore not in conflict with the topological
properties of the field space.

\smallskip
\noindent
(f)~{\itshape Renormalization group}.
By composing the trivializing map $U=\trans_1(V)$
in the Wilson theory
with its inverse at another value of the gauge coupling,
one obtains a group of transformations whose only effect
on the action is a shift of the coupling.
The locality properties of these transformations are
not transparent, however, and could be quite different from the ones
of a Wilsonian ``block spin'' transformation.
